\documentclass[10pt]{article}

\usepackage[margin=0.68in]{geometry}
\usepackage{amsmath,amssymb}
\usepackage{microtype}

\newcommand{\T}{\mathcal{T}}
\newcommand{\C}{\mathbb{C}}
\newcommand{\Hh}{\mathcal{H}}

\setlength{\parindent}{1em}
\setlength{\parskip}{0.15em}
\setlength{\abovedisplayskip}{5pt}
\setlength{\belowdisplayskip}{5pt}

\title{\vspace{-1.6em}Complex-Weighted Operational Theories and Quantum Algorithms}
\author{}
\date{\vspace{-2em}}

\begin{document}
\maketitle
\vspace{-1em}

\paragraph{Complex-weighted operational semantics.}
Suppose that the weighting construction for an operational theory is
specialized to the weight object $\Omega=\C.$
Let $\mathsf{Proc}=\T(1,P)$ be the set of closed processes and let $\Hh$ be
the Hilbert space with orthonormal basis
$\{|p\rangle:p\in\mathsf{Proc}\}$.  A complex-valued weight map on rewrites
induces the operator
\[
A|p\rangle
=
\sum_{r:s(r)=p}\omega(r)|t(r)\rangle.
\]
If several rewrites have the same source and target, their amplitudes add.
Thus the quiver structure of the operational theory becomes a complex-weighted
adjacency operator.  Local summability guarantees that the displayed sum is
well defined on basis vectors, while boundedness of $A$ is a separate
condition needed to extend it to all of $\Hh$.

If $\|A\|\leq 1$, the contraction $A$ has the standard unitary dilation on
$\Hh\oplus\Hh$,
\[
U=
\begin{pmatrix}
A & (I-AA^*)^{1/2}\\
(I-A^*A)^{1/2} & -A^*
\end{pmatrix}.
\]
For an arbitrary bounded $A$, one first replaces it by $A/\alpha$ for
$\alpha\geq\|A\|$.  Identifying $\Hh\oplus\Hh$ with
$\Hh\otimes\C^2$, applying the dilation with a fresh ancillary qubit at each
step and postselecting every ancilla onto $|0\rangle$ realizes
$A^n/\alpha^n$ after $n$ steps.  When the weighted rewrite operator is already
unitary, no dilation or postselection is required.

\section*{Direct encoding of Grover search}

Let $X$ be a finite search type with $N$ closed values, and introduce two
families of process constructors
\[
Q_O,Q_D:X\to P.
\]
The two families distinguish the oracle and diffusion phases of a Grover
iteration.  Let $f:X\to\{0,1\}$ be the search predicate.  The oracle is
represented by the rewrite family
\[
o(x):Q_O(x)\longrightarrow Q_D(x),
\qquad
\omega(o(x))=(-1)^{f(x)}.
\]
Consequently the induced map from the $Q_O$ sector to the $Q_D$ sector is the
usual phase oracle
\[
O_f|x\rangle=(-1)^{f(x)}|x\rangle.
\]

The diffusion operator has an equally direct quiver presentation.  Introduce
a rewrite
\[
d(x,y):Q_D(x)\longrightarrow Q_O(y)
\]
for every pair $x,y\in X$, with
\[
\omega(d(x,y))=\frac{2}{N}-\delta_{xy}.
\]
For a fixed source $x$, there are therefore $N$ outgoing rewrites: amplitude
$2/N$ to every distinct $y$, and amplitude $2/N-1$ on the self-loop.  The
induced operator is
\[
D|x\rangle
=
\sum_{y\in X}
\left(\frac{2}{N}-\delta_{xy}\right)|y\rangle
=
\frac{2}{N}\sum_{y\in X}|y\rangle-|x\rangle.
\]
Writing
\[
|s\rangle=\frac{1}{\sqrt N}\sum_{y\in X}|y\rangle,
\]
we obtain
\[
D=2|s\rangle\langle s|-I,
\]
which is exactly Grover's diffusion operator.

In a pi-calculus-like syntax with internal complex-valued \texttt{where}
expressions, the rule can be presented schematically as
\[
\texttt{x,y:X |- diffuse(x,y): Q\_D(x) \string~> Q\_O(y)
        where 2/N - delta(x,y).}
\]
The significant feature is that $y$ is a finite parameter of the rewrite rule
that is not determined by matching its source.  Its ground instances generate
the complete family of outgoing edges from $Q_D(x)$.  This use of a
target-only parameter is sufficient to express dense linear operators.

On the Hilbert space
$\Hh_O\oplus\Hh_D$ spanned by the two process families, one small operational
step has block form
\[
A=
\begin{pmatrix}
0&D\\
O_f&0
\end{pmatrix}.
\]
Since $D$ and $O_f$ are unitary, so is $A$, and
\[
A^2=
\begin{pmatrix}
DO_f&0\\
0&O_fD
\end{pmatrix}.
\]
Thus, starting in the oracle sector, two small-step reductions implement one
Grover iterate.  After approximately
$\frac{\pi}{4}\sqrt{N/M}$ iterations, where $M$ values satisfy $f$, measurement
has the usual Grover success probability.  The initial uniform state may be
taken as the initial state of the semantics, or prepared by the familiar
Hadamard construction when $N$ is a power of two.

\section*{Encoding Shor period finding}

The same mechanism is expressive enough for Shor's algorithm.  Let
$Q=2^m$ be the size of the exponent register and let $N$ be the integer to be
factored.  For an integer $a$ relatively prime to $N$, use process constructors
whose basis states represent pairs
\[
|x\rangle|z\rangle,
\qquad
0\leq x<Q,\quad 0\leq z<N.
\]
As with Grover, distinct process families can mark successive phases of the
algorithm, ensuring that only the rewrite rules for the current phase are
enabled.

Hadamard rewrite rules prepare
\[
\frac{1}{\sqrt Q}\sum_{x=0}^{Q-1}|x\rangle|1\rangle.
\]
Modular exponentiation can then be represented directly by the deterministic
rewrite family
\[
m_a(x,z):
Q_M(x,z)\longrightarrow
Q_F(x,za^x\bmod N),
\qquad
\omega(m_a(x,z))=1.
\]
Because $a$ is invertible modulo $N$, multiplication by $a^x$ permutes the
second register for each $x$, so this rewrite family induces a unitary
permutation.  Starting with $z=1$, it produces the standard state
\[
\frac{1}{\sqrt Q}
\sum_{x=0}^{Q-1}|x\rangle|a^x\bmod N\rangle.
\]

The inverse quantum Fourier transform on the first register also has a direct
weighted-quiver encoding.  Introduce
\[
q(x,y,z):
Q_F(x,z)\longrightarrow Q_R(y,z)
\]
with
\[
\omega(q(x,y,z))
=
\frac{1}{\sqrt Q}
\exp\left(-\frac{2\pi ixy}{Q}\right).
\]
For each basis process $Q_F(x,z)$ this creates $Q$ outgoing rewrites, and the
induced operator is exactly $F_Q^{-1}\otimes I$.  Measurement of the first
register followed by the usual classical continued-fraction computation
recovers the period with the same analysis as in Shor's algorithm.

This direct Fourier rule establishes expressivity but is exponentially dense
in the number $m$ of exponent qubits.  A succinct operational theory can
instead decompose it into Hadamard and controlled-phase rewrite families,
requiring only polynomially many rule applications.  Modular exponentiation
can likewise be decomposed into reversible controlled modular multiplications.
Thus the operational-theory formalism supports both a direct matrix-level
encoding of Shor's quantum operators and the usual polynomial-size circuit
decomposition.  The role of complex-weighted rewriting is to provide the
linear semantics; efficiency depends, as in the circuit model, on whether the
required rewrite families admit succinct and efficiently realizable
presentations.

\end{document}
